\section{Summary of Observations}

\subsection{Classical computing sample}

\begin{itemize}[leftmargin=1.5em]
    \item \textbf{Case 3 (``named, not handed over'') is the modal outcome for Axis
    A} in this sample: six of the papers that use external data cite it by name only,
    leaving retrieval entirely to the reader.
    \item \textbf{Only two papers achieve Case 2 on Axis A with the strongest form of
    reproducibility} --- a direct, working link. Of these, only one (PDMarks) pins the
    reference to an exact commit hash, guaranteeing a byte-identical copy regardless of
    future changes to the upstream repository.
    \item \textbf{Two papers make an explicit availability claim that does not hold up
    under verification} (Case 5) or resolves only to a conditional promise rather than
    an immediate path (Case 6). Both were only detected as such because the stated
    access paths were independently checked rather than taken at face value ---
    underscoring that a stated ``available at ...'' claim is not equivalent to verified
    availability.
    \item \textbf{Self-generated data (Axis B) is comparatively rare} in this sample
    and, when present, is unevenly documented: one paper links its generator code
    directly (Case 2), one paper describes its generation procedure but links nothing
    (Case 3), and two papers leave their own generated output only conditionally or
    unverifiably accessible (Cases 5 and 6).
    \item \textbf{No paper in this sample returned Case 1 on both axes except the
    purely mathematical optimization paper}, consistent with the expectation that
    Case-1/Case-1 should be reserved for papers with no empirical component at all.
\end{itemize}

\subsection{Quantum computing sample}

\begin{itemize}[leftmargin=1.5em]
    \item \textbf{Axis A is overwhelmingly Case 1 in this sample} --- thirteen of the
    fourteen papers use no external dataset at all; the sole exception (the
    transmon-design paper) names its external source (SQuADDS) without a direct link in
    the reviewed text. This is the sharpest structural difference from the classical
    sample, where external-data use (Cases 2--6 on Axis A) was the norm rather than the
    exception.
    \item \textbf{Axis B (self-generated data) is where nearly all the variation in this
    sample lives}, since almost every paper reports its own newly fabricated,
    measured, or simulated data. Eight of fourteen papers reach Case 2 (directly
    accessible), consistent with a real, if inconsistently applied, community norm of
    depositing experimental data to Zenodo or an institutional repository under a
    formal ``Data availability'' heading.
    \item \textbf{Partial releases (data linked, with additional data gated behind a
    request) appear three times} (charge-parity detection, magnon squeezing, and,
    differently, the CPMG-scaling paper's pending archival deposit) --- distinct from
    the classical sample, where a request-only outcome (Case 6) was typically total
    rather than partial.
    \item \textbf{One paper's stated access path is a literal unassigned placeholder}
    (\texttt{10.5061/dryad.XXXXXXXXX}) rather than a broken but once-valid link, a
    distinct failure mode from the classical sample's 404'd GitHub URL: here the paper
    was evidently submitted before the archival deposit was even created, not merely
    before its retrieval was verified.
    \item \textbf{Self-citation to a lab's own prior open-source tooling} (as opposed to
    releasing new data specific to the paper at hand) appears in this sample (the
    LLM-assisted-experiments paper, citing its own group's HAL/Grapher/QuICK releases)
    and was also observed repeatedly in the broader superconducting-qubit review queue
    beyond this fourteen-paper sample --- a citation style that reads, on first pass, like
    a third-party tool reference and requires checking the citing lab's identity against
    the repository owner to correctly attribute.
\end{itemize}

\subsection{Cross-domain comparison}

Read together, the two samples suggest that the two computing communities structure
their data-availability problem differently rather than simply exhibiting different
overall care. The classical sample's difficulty is overwhelmingly an Axis A problem:
papers routinely build on named, citable, but unlinked external benchmarks and
datasets. The quantum sample's difficulty is instead concentrated in Axis B: since these
papers almost always generate their own new experimental or simulated data, the
practice being tested is not ``did you cite your source'' but ``did you actually deposit
what you produced,'' and the sample shows a genuine, if partial and unevenly applied,
norm of doing so (Case 2 on Axis B in 8/14 papers). A same-shape comparison across the
full review queues, once both are fully verified, is needed before this observation can
be treated as more than suggestive --- both samples reported here are small,
non-random, hand-selected subsets chosen for outcome diversity rather than
representativeness (see Limitations).
